% ---- Document-specific additions to the shared preamble ---------------------
\hypersetup{
    pdftitle    = {latex-twocol-cjk: A Claude Skill for IEEE-Style Two-Column CJK-Latin Documents},
    pdfauthor   = {markgoob},
    pdfsubject  = {Introduction to the latex-twocol-cjk skill},
    pdfkeywords = {Claude skill, LuaLaTeX, babel, fontspec, CJK, Traditional Chinese},
}

% Headings are mostly Chinese: \bfseries instead of \scshape / \itshape (SKILL.md §4).
\titleformat{\section}[block]
  {\centering\normalsize\bfseries}{\thesection.}{0.5em}{}
\titleformat{\subsection}[block]
  {\normalsize\bfseries}{\Alph{subsection}.}{0.5em}{}

% Extra TikZ libraries for the diagrams in this document.
\usetikzlibrary{shapes.geometric, decorations.pathreplacing}

\title{\textbf{latex-twocol-cjk: A Claude Skill for IEEE-Style\\
Two-Column CJK--Latin Documents\\[2pt]
latex-twocol-cjk：中英混排雙欄技術文件的 Claude skill}}
\author{
  markgoob\\
  latex-twocol-cjk 專案\\
  \texttt{https://github.com/markgoob/latex-twocol-cjk}
}
\date{}

% ============================== END PREAMBLE =================================
\begin{document}
\raggedbottom

\twocolumn[
  \begin{@twocolumnfalse}
    \maketitle
    \vspace{-1.5em}

    \begin{abstract}
    latex-twocol-cjk is a Claude skill for A4, two-column, IEEE-style documents that mix Traditional Chinese with English: internal whitepapers, technical reports, course papers. It packages a LuaLaTeX preamble that compiles with zero warnings under both Noto font naming conventions, a troubleshooting file that maps 24 real log messages to causes and fixes, and four scripts for font checking, containerised building, log gating and packaging. This note introduces the skill: what is in the package, the three-layer structure of a document built from it, the design decisions in the preamble and the failure each one prevents, the source-formatting and prose rules specific to Chinese text, and how every change is verified by a log gate that fails on any warning rather than on exit code alone. It closes with the known limits: the output is IEEE-styled, not IEEE-submittable. This PDF was produced with the skill and passes the same gate.
    \end{abstract}

    \begin{IEEEkeywords}
    Claude skill, LuaLaTeX, babel, fontspec, CJK typesetting, Traditional Chinese, reproducible builds
    \end{IEEEkeywords}

    \begin{cnabstract}
    latex-twocol-cjk 是一個 Claude skill，用來寫 A4 雙欄、IEEE 風格、中文夾英文的技術文件：內部白皮書、技術報告、課堂論文。套件包含一份在兩套 Noto 字型命名慣例下都以零警告編譯的 LuaLaTeX preamble、一份把 24 條真實 log 訊息對應到成因與修法的錯誤對照表，以及檢查字型、容器內編譯、log 閘門、打包用的四支腳本。本文介紹這個 skill：套件裡有什麼、用它寫出來的文件分成哪三層、preamble 裡每個設計決策各擋掉哪一種失敗、中文原始碼與行文的專屬規則，以及每次改動如何用「任何警告即失敗」的 log 閘門驗證，而不是只看 exit code。最後列出已知限制：產出是 IEEE 風格，不是可投稿的格式。這份 PDF 本身就由這個 skill 產出，並通過同一道閘門。
    \end{cnabstract}

    \begin{cnkeywords}
    Claude skill、LuaLaTeX、babel、fontspec、中日韓排版、繁體中文、可重現建置
    \end{cnkeywords}
  \end{@twocolumnfalse}
]

% =============================================================================
\section{這個 skill 是什麼}
\label{sec:what}

Claude 的 skill 是一組寫在 \texttt{SKILL.md} 裡的工作指引，加上它引用的參考檔與腳本。使用者提出符合描述的需求時，模型先載入這組指引再動手。latex-twocol-cjk 針對的是一種很具體的文件：A4 雙欄、IEEE 風格、中文內文夾英文術語的技術報告、白皮書與課堂論文。它把一份驗證過的 LuaLaTeX 範本、一份錯誤對照表和四支腳本包在一起，讓模型寫這類文件時不必每次重新摸索字型與版面。

這個 skill 有一段該交代的來歷。第一版建在 XeLaTeX 加 \texttt{bidi=basic} 的設定上，而那個組合根本產不出 PDF；第二版改用 LuaLaTeX 重寫，才有了現在這份範本。之後的每一次改動都要通過同一套 log 閘門，範本自己在兩套字型命名慣例下都維持零警告。目前版本是 v1.1.0，以 MIT 授權公開~\cite{repo}，並附繁體中文說明~\cite{repo-zh}。

適用範圍要先講清楚。範本建在 \texttt{article} 上而不是 IEEEtran，產出看起來像 IEEE 論文，適合內部報告與課堂論文；它不是可以投稿的 camera-ready 格式，邊界、摘要寬度、圖說度量都與投稿規格不同。真要投稿，請用 IEEEtran 並只移植字型那一段。

\section{套件內容}
\label{sec:contents}

套件只有八個檔案（Table~\ref{tab:files}）。\texttt{SKILL.md} 是規則本體，九個編號章節按寫文件的實際順序排列：引擎與語系、字型、標題區、標題與交叉引用、浮動體、圖與表、程式碼與數學與 TikZ、中文原始碼的兩條規則、中文行文，再加一節「不改範本就擴充 preamble」與一節編譯驗證。每條規則都附理由，因為多數規則單看會覺得多餘，知道它擋掉哪個失敗才會照做。

\begin{table}[!t]
\centering
\caption{套件內的檔案與用途}
\label{tab:files}
\small
\begin{tabularx}{\columnwidth}{@{}l Y@{}}
\toprule
\textbf{檔案} & \textbf{用途} \\
\midrule
\texttt{SKILL.md} & 規則本體，按寫文件的順序排列 \\
\texttt{references/template.tex} & 完整且驗證過的文件；preamble 直接複製 \\
\texttt{references/troubleshooting.md} & 24 條錯誤：log 原文、成因、修法 \\
\texttt{scripts/check-fonts.ps1}、\texttt{.sh} & 列出候選字型家族裝了哪些 \\
\texttt{scripts/build.ps1} & 沒裝 TeX 的機器用的容器化編譯 \\
\texttt{scripts/log-gate.sh} & 對 log 做五項計數，任一非零即失敗 \\
\texttt{scripts/package-skill.ps1} & 打包成安裝程式接受的 \texttt{.skill} 檔 \\
\bottomrule
\end{tabularx}
\end{table}

\texttt{template.tex} 不只是 preamble，它本身是一份完整可編譯的兩頁文件，示範了範本支援的每一種元素：中英雙摘要、TikZ 架構圖、子圖、siunitx 對齊的數值表、帶中文註解的 Verbatim、方程式、中文參考文獻。不用 Claude 也能拿它當一般 LaTeX 範本。

\texttt{troubleshooting.md} 的每一條都從實際出現過的 log 訊息出發，例如 \texttt{Missing character}、\texttt{Font shape undefined}、babel 的 bidi 錯誤、fontspec 的 font cannot be found，再說成因與修法。它的價值在於「先搜 log 原文」這個入口：使用者看到的永遠是訊息，不是成因。

\section{使用流程}
\label{sec:workflow}

用這個 skill 寫出來的文件分成三層（Fig.~\ref{fig:layers}）。第一層是範本的 preamble，從檔案開頭複製到 \texttt{TITLE \& AUTHORS} 標記為止，一個字都不改。第二層是文件自己的設定：第二個 \verb|\hypersetup| 蓋掉佔位的 PDF metadata，需要的話再覆寫標題格式或補 TikZ 函式庫。第三層才是內文。

\begin{figure}[!t]
\centering
\resizebox{0.92\columnwidth}{!}{%
\begin{tikzpicture}[
  layer/.style={draw=blue!50!cyan!60, fill=blue!6, rounded corners=4pt,
    minimum width=5.2cm, minimum height=1.05cm, font=\sffamily\small,
    text=blue!60!black, line width=0.6pt, align=center},
  own/.style={layer, draw=green!50!black!60, fill=green!6, text=green!40!black},
  note/.style={font=\sffamily\scriptsize, text=black!60, align=left, anchor=west}]
\node[layer] (p) {範本 preamble\\[-1pt]{\scriptsize 複製到 TITLE \& AUTHORS 標記為止}};
\node[own, below=0.3cm of p] (o) {文件自己的設定\\[-1pt]{\scriptsize \textbackslash hypersetup、\textbackslash title、\textbackslash titleformat 覆寫}};
\node[own, below=0.3cm of o] (b) {內文\\[-1pt]{\scriptsize 標題區、章節、圖表、參考文獻}};
\node[note, right=0.35cm of p] {逐字複製，\\不修改};
\node[note, right=0.35cm of o] {後定義者勝，\\覆寫佔位值};
\node[note, right=0.35cm of b] {中文段落\\在原始碼裡單行};
\draw[decorate, decoration={brace, amplitude=5pt, mirror}, black!40, line width=0.6pt]
  ([xshift=-0.3cm]o.north west) -- ([xshift=-0.3cm]b.south west)
  node[midway, xshift=-0.55cm, font=\sffamily\scriptsize, text=black!60, rotate=90, anchor=south] {你寫的部分};
\end{tikzpicture}%
}
\caption{用 skill 寫出的文件分三層。只有第一層來自範本，而且是原封不動地複製；後兩層屬於文件本身。}
\label{fig:layers}
\end{figure}

骨架長這樣：
\begin{Verbatim}
% 1. 範本 preamble，複製到 TITLE & AUTHORS 標記為止
% 2. 自己的 \hypersetup、\title、\author、\date{}
\begin{document}
\raggedbottom
\twocolumn[ ... 標題、摘要、關鍵詞 ... ]
% 3. 內文
\end{document}
\end{Verbatim}

編譯之前先跑字型檢查，因為缺一個家族是 fatal error 而不是警告，早發現省時間。之後用 latexmk 編譯；本機沒裝 TeX 的話，\texttt{build.ps1} 會在 TeX Live 容器內編譯並掛載主機的字型目錄，看到的是使用者實際擁有的字型環境（Fig.~\ref{fig:flow}）。
\begin{Verbatim}
scripts\check-fonts.ps1                       # Windows
latexmk -lualatex -interaction=nonstopmode main.tex
scripts\build.ps1 main.tex                    # 沒裝 TeX 時
\end{Verbatim}
Overleaf 會忽略 \texttt{\% !TeX program} 那行，得手動把編譯器切成 LuaLaTeX；它的 TeX Live 已內建 Noto CJK 字型。

\begin{figure*}[!t]
\centering
\begin{subfigure}[t]{0.47\textwidth}
  \centering
  \resizebox{\linewidth}{!}{%
  \begin{tikzpicture}[node distance=0.4cm,
    stage/.style={draw=blue!50!cyan!60, fill=blue!6, rounded corners=4pt,
      minimum width=2.5cm, minimum height=0.7cm, font=\sffamily\small,
      text=blue!60!black, line width=0.6pt, align=center},
    io/.style={stage, draw=black!40, fill=white, text=black!80},
    arr/.style={->, >=Stealth, black!50, line width=0.6pt}]
  \node[io] (a) {check-fonts};
  \node[stage, right=of a] (b) {寫 body\\[-1pt]{\scriptsize 中文段落單行}};
  \node[stage, right=of b] (c) {latexmk\\[-1pt]{\scriptsize 或 build.ps1}};
  \node[io, right=of c] (d) {PDF};
  \draw[arr] (a) -- (b); \draw[arr] (b) -- (c); \draw[arr] (c) -- (d);
  \end{tikzpicture}}
  \caption{撰寫與編譯}
  \label{fig:flow-a}
\end{subfigure}\hfill
\begin{subfigure}[t]{0.47\textwidth}
  \centering
  \resizebox{\linewidth}{!}{%
  \begin{tikzpicture}[node distance=0.5cm,
    stage/.style={draw=blue!50!cyan!60, fill=blue!6, rounded corners=4pt,
      minimum width=2.6cm, minimum height=0.7cm, font=\sffamily\small,
      text=blue!60!black, line width=0.6pt, align=center},
    dec/.style={diamond, aspect=2.2, draw=red!50, fill=red!5, font=\sffamily\small,
      text=red!60!black, line width=0.6pt, inner sep=1pt},
    ok/.style={stage, draw=green!50!black!60, fill=green!6, text=green!40!black},
    arr/.style={->, >=Stealth, black!50, line width=0.6pt},
    lbl/.style={font=\sffamily\scriptsize, text=black!55}]
  \node[stage] (c) {編譯};
  \node[stage, right=of c] (g) {grep log\\[-1pt]{\scriptsize 五項計數}};
  \node[dec, right=of g] (q) {全為零？};
  \node[ok, right=of q] (y) {交付};
  \node[stage, below=0.55cm of q] (n) {修正};
  \draw[arr] (c) -- (g); \draw[arr] (g) -- (q);
  \draw[arr] (q) -- node[lbl, above] {是} (y);
  \draw[arr] (q) -- node[lbl, right] {否} (n);
  \draw[arr] (n) -| (c);
  \end{tikzpicture}}
  \caption{log 閘門}
  \label{fig:flow-b}
\end{subfigure}
\caption{工作流程。(a) 先查字型再寫再編；(b) 編譯後不看 exit code，而是對 log 做五項計數，任何一項非零就回頭修。}
\label{fig:flow}
\end{figure*}

\section{範本裡的設計決策}
\label{sec:design}

preamble 裡每一個不尋常的設定都對應一個實際發生過的失敗（Table~\ref{tab:design}）。這一節挑最要緊的幾個說明。

\begin{table}[!t]
\centering
\caption{設計決策與它擋掉的失敗}
\label{tab:design}
\small
\begin{tabularx}{\columnwidth}{@{}Y Y@{}}
\toprule
\textbf{決策} & \textbf{擋掉的失敗} \\
\midrule
LuaLaTeX，不用 XeLaTeX & \texttt{bidi=basic} 的硬錯誤；\texttt{onchar} 安靜失效 \\
\texttt{english} 主語言，中文為次語言 & 圖說變成 図、表、参考文献，日期成日式 \\
\texttt{\textbackslash PickFont} 候選鏈 & 兩套 Noto 命名造成 font cannot be found \\
釘住 \texttt{wght} 軸 & 可變字型預設實例是 ExtraLight \\
rm、sf、tt 三家族雙語系宣告 & \texttt{\textbackslash texttt}、Verbatim、TikZ 裡的中文消失 \\
luaotfload 字形後備 & 區域子集缺字，例如日文的 図 \\
caption 套件設定 & article 預設的「Figure 1:」 \\
重定義 \texttt{\textbackslash thesection} & 標題印 II，內文印 Section 2 \\
\texttt{\textbackslash raggedbottom} 與浮動體參數 & 浮動體周圍的大片空白 \\
\texttt{intraspace=\{0 .22 .10\}} & 中文行的 badness 10000 underfull \\
\bottomrule
\end{tabularx}
\end{table}

\subsection{引擎與語系}

babel 的 \texttt{onchar=ids fonts} 是整套設定的核心：遇到漢字切中文字型，遇到拉丁字母切回西文字型。它只有 LuaTeX 支援~\cite{babel}，在 XeLaTeX 下不報錯但從未生效；同一份常見設定裡的 \texttt{bidi=basic} 更是直接讓 babel 停機。範本因此指定 LuaLaTeX，並且不傳任何 \texttt{bidi=} 選項。

主語言是 \texttt{english}，中文以 \texttt{chinese-traditional} 掛成次語言。這樣圖表標題與日期維持英文（IEEE 要的），漢字則拿到正確字型、斷行規則與 PDF 裡的 zh-Hant 標記。babel 內建 zh-Hant 與 zh-Hans 兩個語系，沒有理由拿 \texttt{japanese} 充數。

\subsection{字型}

字型是這個 skill 花最多篇幅的地方，因為它有三層坑。第一層是命名：同一套 Noto CJK 字型，noto-cjk 專案與 Linux 套件叫 Noto Serif CJK TC，Google Fonts 叫 Noto Serif TC，後者才是 Windows 使用者裝到的~\cite{notocjk}。範本用 \texttt{\textbackslash PickFont} 逐一探測候選家族，一路退到 Windows 內建的新細明體與微軟正黑體。

第二層是權重。Google Fonts 版是單一可變字型檔，第一個具名實例是 ExtraLight；照家族名查找就會命中它，中文內文比旁邊的拉丁文字明顯細，log 完全沒有警告。範本把權重軸釘在 400，粗體釘在 700，順帶換到真正的粗體。靜態字型會忽略這個設定，所以同一行寫法對兩類字型都安全。

第三層是覆蓋範圍。區域子集沒有日文新字體與許多簡體字形，缺的字會從 PDF 消失。v1.1.0 起範本會偵測覆蓋較廣的第二字型並以 luaotfload~\cite{luaotfload} 掛成字形後備；一個都沒有時只寫一行 Info，行為與之前相同。另外，三個家族都要對兩種文字系統宣告，否則 \texttt{\textbackslash texttt}、Verbatim 與 TikZ 標籤裡的中文會落到沒有漢字的西文字型。

\subsection{IEEE 細節}

IEEE 的圖說是「Fig. 1.」而不是 article 預設的「Figure 1:」，表格是置中的「TABLE I」加小型大寫標題，這些需要 caption 套件明確設定。\texttt{\textbackslash thesection} 要重新定義成羅馬數字，否則 titlesec 只改了標題顯示，內文的交叉引用仍印阿拉伯數字。雙欄版面另有一組浮動體參數與 \texttt{\textbackslash raggedbottom}，欄距設 16 pt，並載入 microtype 減少窄欄的溢出。

\FloatBarrier
\section{中文行文}
\label{sec:prose}

中文原始碼有兩條規則，都來自實際寫過的長文件。第一，不要在兩個漢字之間換行：TeX 在 tokenise 階段就把行尾變成空格，比 babel 看到文字更早，每個被編輯器折行的中文段落都會在 PDF 上留下句子中間的空格。本文的每個中文段落在原始碼裡都是一整行。第二，中英交界要自己打半形空格，babel 不像 xeCJK 會自動插入；這也正是第一條規則不能套用的地方。行距鬆散時該調的是語系的 \texttt{intraspace}，不是 \verb|\sloppy|。

用字與標點的規範寫在 \texttt{SKILL.md} 第 9 節：臺灣用語而非大陸用語，中文句子裡用全形標點，半形標點只留給英文與英文之間。兩處是 LaTeX 特有的陷阱：反引號式的引號輸入會排成拉丁引號而不是「」，\verb|\ldots| 會排成拉丁省略號而不是⋯⋯。摘要最容易招來空話，第 9 節直接列出該刪的句型。如果同時裝了 be-human-v1~\cite{behuman}，散文的部分預設交給它，這個 skill 只保留影響版面的判斷；兩者互不相依。

% \balance goes in the first column of the last page (SKILL.md §10); §VI is
% where that column begins in the current layout. If the layout shifts and this
% lands in the second column, the log gate catches the resulting warning.
\balance
\section{驗證與維護}
\label{sec:verify}

exit code 為 0 不代表文件正確。LuaLaTeX 找不到字型時照樣成功結束，只是把字丟掉。因此驗證的單位是 log：對五個模式計數，任何一項非零就算失敗。
\begin{Verbatim}
for p in 'Missing character' 'Font shape .* undefined' \
         'Overfull' 'Underfull'; do
  grep -c "$p" main.log
done
grep -cE '^(LaTeX|Package|Class) .*Warning' main.log
\end{Verbatim}
\texttt{build.ps1} 編完會直接印出這五個數字。GitHub Actions 在每次改動範本時跑同一套檢查，用的是 Debian 的 \texttt{fonts-noto-cjk}，也就是 Noto Serif CJK TC 那套命名；本機 Windows 用的是 Google Fonts 命名。兩邊各走候選鏈的一半（Table~\ref{tab:resolve}），所以兩條路徑都持續受測，一次 CI 約 1 分 45 秒。

\begin{table}[!t]
\centering
\caption{兩個環境下各字型槽的實際解析結果}
\label{tab:resolve}
\small
\setlength{\tabcolsep}{3pt}
\begin{tabularx}{\columnwidth}{@{}l Y Y@{}}
\toprule
\textbf{槽} & \textbf{Windows 11}\rlap{$^{\star}$} & \textbf{Debian CI} \\
\midrule
拉丁 rm & Noto Serif & Noto Serif \\
拉丁 sf & Noto Sans & Noto Sans \\
拉丁 tt & Noto Sans Mono & Noto Sans Mono \\
中文 rm & Noto Serif TC & Noto Serif CJK TC \\
中文 sf & Noto Sans TC & Noto Sans CJK TC \\
中文 tt & Noto Sans TC & Noto Sans Mono CJK TC \\
字形後備 & Microsoft JhengHei & Noto Sans CJK TC \\
\bottomrule
\end{tabularx}

\vspace{4pt}
\raggedright
{\scriptsize $^{\star}$\,經 \texttt{build.ps1} 的容器加掛主機字型的結果。直接在 Windows 上跑 \texttt{check-fonts.ps1} 時，拉丁襯線會退到 Times New Roman，因為 Noto Serif 來自 TeX Live 的 noto 套件而非系統。}
\end{table}

打包也有一個坑。Windows PowerShell 5.1 的 \texttt{Compress-Archive} 會把反斜線寫進壓縮檔的路徑，安裝程式看到的是 Zip file contains path with invalid characters。\texttt{package-skill.ps1} 逐一手動建立條目並強制正斜線，寫檔前自我驗證：有反斜線、絕對路徑、目錄穿越，或根目錄缺 \texttt{SKILL.md}，就刪檔中止。Release 上的資產都由它產生，並另以 Linux 的 unzip 驗過。

\section{限制}
\label{sec:limits}

有幾件事這個 skill 明確不做，或還沒做。產出不是可投稿的 IEEE 格式，這點前面說過。字形後備在粗體文字裡會以第二字型的常規字重印出。韓文需要另配 KR 字型，範本只給了語系名稱。書目只示範 thebibliography，混合語言書目該用的 biblatex 加 biber 路線只在規則裡提了一句。草稿模式（行號、浮水印）與 TikZ 外部化也還沒寫進去。混排標題若保留 \verb|\scshape|，英文半邊會變小型大寫而中文半邊維持原樣，範本的建議是換成 \verb|\bfseries|，本文照做了。

\section{結語}
\label{sec:end}

這個 skill 的價值不在任何單一規則，而在每條規則都對應一個重現過的失敗，並且用同一道閘門守著。這份 PDF 以範本的 preamble 加上本文內容，在同一個容器內以 latexmk 編譯，五項計數全為零；它示範的每一個元素都是 skill 支援範圍內的東西。

% --- BIBLIOGRAPHY: numbered in order of first citation ---
% \balance decided after inspecting the unbalanced output (SKILL.md §10).
\begingroup
\raggedright
\small
\begin{thebibliography}{9}
\bibitem{repo}
markgoob, ``latex-twocol-cjk,'' v1.1.0, GitHub, 2026. [Online]. Available:
\url{https://github.com/markgoob/latex-twocol-cjk}. [Accessed: Sep.~3, 2026].

\bibitem{repo-zh}
markgoob，「latex-twocol-cjk 繁體中文說明」，GitHub，2026。[Online]. Available:
\url{https://github.com/markgoob/latex-twocol-cjk/blob/main/README.zh-TW.md}.
[Accessed: Sep.~3, 2026] (in Chinese).

\bibitem{babel}
J. Bezos and J. L. Braams, \textit{Babel: Localization and
internationalization}, v26.9, CTAN, 2026. [Online]. Available:
\url{https://ctan.org/pkg/babel}. [Accessed: Sep.~3, 2026].

\bibitem{notocjk}
Google, ``Noto CJK,'' GitHub repository. [Online]. Available:
\url{https://github.com/notofonts/noto-cjk}. [Accessed: Sep.~3, 2026].

\bibitem{luaotfload}
The LuaLaTeX Development Team, \textit{The luaotfload package}, CTAN.
[Online]. Available: \url{https://ctan.org/pkg/luaotfload}.
[Accessed: Sep.~3, 2026].

\bibitem{behuman}
markgoob, ``be-human-v1,'' GitHub, 2026. [Online]. Available:
\url{https://github.com/markgoob/be-human-v1}. [Accessed: Sep.~3, 2026].
\end{thebibliography}
\endgroup

\end{document}
